\section{矩、弱收敛与测度识别}\label{sec:07-03}

一个有限测度可以在许多不同的测试函数族上积分：多项式给出矩，复指数给出 Fourier 系数，半直线指标给出
分布函数。测试函数族太小时，不同测度可能无法区分；即使测试函数族足够大，还要回答这些数值是否真的来自
某个测度。本节把这两个问题分开处理。第一小节研究矩数据的正性与表示，第二小节研究测度列的选择
与弱收敛，最后一小节比较几类变换的识别能力，并用 Carleman 条件修补无界支集上的矩不唯一性。

Hamburger、Stieltjes、Hausdorff 与 Herglotz 的矩问题分别把支集限制在
$\R$、$[0,\infty)$、$[0,1]$ 与圆周 $\T$。它们都可写成“正线性泛函是否来自积分”，但支集不同，
可检验的正性条件也不同。尤其不能把无界支集问题中的存在性、质量控制与唯一性压缩成同一句表示
定理。

\subsection{矩问题、正性与 Hankel/Toeplitz 结构}\label{subsec:7-3-1}
\index{矩问题}\index{Hankel 矩阵}\index{Toeplitz 矩阵}\index{Herglotz 表示}

先从最小的可计算模型开始。若 $\mu$ 是 $\R$ 上的有限正测度并且所需矩有限，则任意多项式
$p$ 都满足 $\int |p|^2\,\dd\mu\geq0$。展开平方后，这个不等式只涉及有限多个矩；因此每次检验
都是有限维矩阵的半正定性。圆周上的三角多项式给出同样的机制，只是矩阵由 Hankel 型变成
Toeplitz 型。

\begin{example}[Dirac 矩]\label{ex:7-3-1-1}
对 $\mu=\delta_a$，
\[
 c_n=\int_\R x^n\,\delta_a(\dd x)=a^n.
\]
令 $v_N(a)=(1,a,\ldots,a^N)^{\mathsf T}$，则
\[
 H_N=(c_{j+k})_{0\leq j,k\leq N}
     =v_N(a)v_N(a)^{\mathsf T}.
\]
所以 $H_N$ 半正定且秩为一。更一般地，若
$\mu=\sum_{\ell=1}^m w_\ell\delta_{a_\ell}$，其中 $w_\ell>0$ 且 $a_\ell$ 两两不同，则
\[
 H_N=V_N\operatorname{diag}(w_1,\ldots,w_m)V_N^{\mathsf T},
 \qquad (V_N)_{j\ell}=a_\ell^j.
\]
Vandermonde 行列式说明 $N\geq m-1$ 时 $V_N$ 的列线性无关，因而
$\operatorname{rank}H_N=m$；对任意 $N$，秩至多为 $m$。矩阵的有限秩正是有限原子数留下的代数
痕迹。
\end{example}

\begin{example}[圆周正定序列]\label{ex:7-3-1-2}
设 $\mu$ 是 $\T=\R/(2\pi\Z)$ 上的有限正测度，并令
\[
 c_n=\int_\T e^{-int}\,\mu(\dd t),\qquad n\in\Z.
\]
于是 $c_{-n}=\overline{c_n}$。对任意 $a_0,\ldots,a_N\in\C$，直接展开得
\begin{align*}
 \sum_{j,k=0}^N\overline{a_j}c_{j-k}a_k
 &=\int_\T\sum_{j,k=0}^N
   \overline{a_j}a_k e^{-i(j-k)t}\,\mu(\dd t)\\
 &=\int_\T\left|\sum_{k=0}^N a_ke^{ikt}\right|^2\mu(\dd t)\geq0.
\end{align*}
下文的 Herglotz 定理说明，这些有限二次型的正性不仅必要，而且足以构造出唯一的 $\mu$。
\end{example}

\begin{example}[对数正态不定矩]\label{ex:7-3-1-3}
令 $Y\sim N(0,1)$、$X=e^Y$。$X$ 的密度为
\[
 f(x)=\frac{1}{x\sqrt{2\pi}}
      \exp\!\left(-\frac{(\log x)^2}{2}\right),\qquad x>0,
\]
且
\[
 \int_0^\infty x^n f(x)\,\dd x=\E e^{nY}=e^{n^2/2},
 \qquad n=0,1,2,\ldots.
\]
对 $|\varepsilon|\leq1$ 定义 Stieltjes 族
\[
 f_\varepsilon(x)=f(x)\bigl[1+\varepsilon\sin(2\pi\log x)\bigr].
\]
括号内非负。对每个整数 $n\geq0$，扰动项的 $n$ 阶矩为
\begin{align*}
 \int_0^\infty x^nf(x)\sin(2\pi\log x)\,\dd x
 &=\Im\E e^{(n+2\pi i)Y}\\
 &=\Im\exp\!\left(\frac{(n+2\pi i)^2}{2}\right)\\
 &=e^{(n^2-4\pi^2)/2}\sin(2\pi n)=0.
\end{align*}
取 $n=0$ 还说明 $f_\varepsilon$ 的积分为一。因此所有 $f_\varepsilon$ 都是概率密度，并共享矩
$e^{n^2/2}$；当 $\varepsilon\neq\eta$ 时，两密度在一个正 Lebesgue 测度集合上不同。这个计算表明：
即使每一阶矩都存在且全部相同，分布律仍可能不同。
\end{example}

\begin{definition}[矩泛函]\label{def:7-3-1-1}
给定实数列 $(c_n)_{n\geq0}$，在 $\C[x]$ 上定义
\[
 L\!\left(\sum_{k=0}^N a_kx^k\right)=\sum_{k=0}^Na_kc_k.
\]
对多项式 $p(x)=\sum a_kx^k$，记
$\overline p(x)=\sum\overline{a_k}x^k$。若
\[
 L(\overline p\,p)\geq0\qquad(p\in\C[x]),
\]
则称 $L$ 为\emph{正矩泛函}。
\end{definition}

\begin{definition}[Hankel 与 Toeplitz 矩阵]\label{def:7-3-1-2}
单边矩序列 $(c_n)_{n\geq0}$ 的第 $N$ 个 \emph{Hankel 矩阵}是
\[
 H_N=(c_{j+k})_{0\leq j,k\leq N}.
\]
满足 $c_{-n}=\overline{c_n}$ 的双边序列的第 $N$ 个 \emph{Toeplitz 矩阵}是
\[
 T_N=(c_{j-k})_{0\leq j,k\leq N}.
\]
所谓所有截断半正定，是指 $a^*H_Na\geq0$ 或 $a^*T_Na\geq0$ 对每个 $N$ 与每个相应有限向量
$a$ 成立。
\end{definition}

\begin{definition}[确定与不确定矩问题]\label{def:7-3-1-3}
若一个矩序列存在满足指定支集条件的正表示测度，并且该测度唯一，则称相应矩问题是
\emph{确定的}（determinate）；若存在至少两个不同的表示测度，则称为
\emph{不确定的}（indeterminate）。没有表示测度的序列不归入这两个名称。存在性与确定性是两个独立
问题：先要证明表示存在，随后才谈得上它是否唯一。
\end{definition}

\begin{proposition}[Hankel 正性必要性]\label{proposition:7-3-1-1}
设 $\mu$ 是 $\R$ 上的正测度，且 $c_n=\int x^n\,\dd\mu$ 对所有 $n\geq0$ 有限。则每个
$H_N$ 半正定。若进一步 $\supp\mu\subset[0,\infty)$，则每个移位 Hankel 矩阵
\[
 H_N^{(1)}=(c_{j+k+1})_{0\leq j,k\leq N}
\]
也半正定。
\end{proposition}

\begin{proof}
对 $p(x)=\sum_{k=0}^Na_kx^k$，有限求和可直接移入积分，故
\begin{align*}
 a^*H_Na
 &=\sum_{j,k=0}^N\overline{a_j}a_kc_{j+k}
   =\int_\R\sum_{j,k=0}^N\overline{a_j}a_kx^{j+k}\,\mu(\dd x)\\
 &=\int_\R|p(x)|^2\,\mu(\dd x)\geq0.
\end{align*}
若 $\mu$ 支持在 $[0,\infty)$，同样的展开给出
\[
 a^*H_N^{(1)}a=\int_{[0,\infty)}x|p(x)|^2\,\mu(\dd x)\geq0.
\]
两项结论分别成立。
\end{proof}

这里必须保留支集条件。经典 Hamburger 判据以所有 $H_N\succeq0$ 刻画 $\R$ 上的表示；经典
Stieltjes 判据还要求所有 $H_N^{(1)}\succeq0$，以刻画 $[0,\infty)$ 上的表示。完整的充分性证明
需要无界对称算子的方法。圆周与闭区间情形中支集是紧的，因而可以直接用有限正测度和紧性完成构造。

\begin{theorem}[Herglotz 表示]\label{theorem:7-3-1-2}
设双边序列 $(c_n)_{n\in\Z}$ 满足 $c_{-n}=\overline{c_n}$。存在唯一的 $\T$ 上有限正测度
$\mu$ 使
\[
 c_n=\int_\T e^{-int}\,\mu(\dd t),\qquad n\in\Z,
\]
当且仅当所有 Toeplitz 截断 $T_N$ 半正定。
\end{theorem}

\begin{proof}
必要性已在例~\ref{ex:7-3-1-2} 的有限平方展开中证明。以下证明充分性。

先分离零质量情形。对 $n\geq1$，$T_n$ 的第 $0,n$ 行列所成主子式是
\[
 \begin{pmatrix}c_0&c_{-n}\\c_n&c_0\end{pmatrix}.
\]
半正定性给 $c_0\geq0$ 以及
$c_0^2-|c_n|^2\geq0$。若 $c_0=0$，则 $c_n=0$ 对所有 $n\geq1$，再由共轭对称性可知整个序列
为零；零测度给出表示。以下设 $c_0>0$。

对 $N\geq0$ 定义三角多项式
\begin{equation}\label{eq:7-3-herglotz-fejer-density}
 P_N(t)=\frac{1}{N+1}\sum_{j,k=0}^Nc_{j-k}e^{i(j-k)t}.
\end{equation}
把 $a_k=e^{-ikt}$ 代入 $T_N$ 的二次型，得到
$P_N(t)=(N+1)^{-1}a^*T_Na\geq0$。于是
\[
 \mu_N(\dd t)=P_N(t)\frac{\dd t}{2\pi}
\]
是有限正测度。按差 $m=j-k$ 收集式~\eqref{eq:7-3-herglotz-fejer-density} 中的项，有
\[
 P_N(t)=\sum_{|m|\leq N}
 \left(1-\frac{|m|}{N+1}\right)c_me^{imt}.
\]
圆周指数函数的正交性因而给出
\begin{equation}\label{eq:7-3-herglotz-fejer-coefficients}
 \int_\T e^{-int}\,\mu_N(\dd t)
 =
 \begin{cases}
 \left(1-\dfrac{|n|}{N+1}\right)c_n,& |n|\leq N,\\[4pt]
 0,& |n|>N.
 \end{cases}
\end{equation}
特别地，$\mu_N(\T)=c_0$。概率测度 $c_0^{-1}\mu_N$ 全部支持在紧空间 $\T$ 上，故由
Prokhorov 定理~\ref{theorem:5-3-3-2} 可取子列 $N_r\to\infty$，使
$c_0^{-1}\mu_{N_r}$ 弱收敛于某个概率 $\lambda$。令 $\mu=c_0\lambda$。对每个固定 $n\in\Z$，
函数 $t\mapsto e^{-int}$ 有界连续；弱收敛与
式~\eqref{eq:7-3-herglotz-fejer-coefficients} 给出
\[
 \int_\T e^{-int}\,\mu(\dd t)
 =\lim_{r\to\infty}\int_\T e^{-int}\,\mu_{N_r}(\dd t)=c_n.
\]
这就得到表示。

若 $\mu$ 与 $\nu$ 都表示 $(c_n)$，则有限复测度 $\mu-\nu$ 的全部圆周 Fourier 系数为零。
圆周 Fourier 唯一性命题~\ref{proposition:7-2-1-4} 给出 $\mu-\nu=0$，所以表示唯一。
\end{proof}

\begin{theorem}[Hausdorff 矩判据]\label{theorem:7-3-1-3}
实数列 $(c_n)_{n\geq0}$ 是 $[0,1]$ 上某个有限正测度的矩序列，当且仅当
\begin{equation}\label{eq:7-3-hausdorff-complete-monotonicity}
 (-1)^k\Delta^kc_n\geq0\qquad(n,k\geq0),
 \qquad \Delta c_n=c_{n+1}-c_n.
\end{equation}
表示测度唯一，其总质量为 $c_0$。
\end{theorem}

\begin{proof}
若 $c_n=\int_0^1x^n\,\mu(\dd x)$，则有限差分可以逐项移入积分。由
$\Delta^kx^n=x^n(x-1)^k$ 得
\[
 (-1)^k\Delta^kc_n
 =\int_0^1x^n(1-x)^k\,\mu(\dd x)\geq0.
\]
这证明必要性。

反过来，假设式~\eqref{eq:7-3-hausdorff-complete-monotonicity} 成立。若 $c_0=0$，则
$c_n\geq0$，而 $c_n-c_{n+1}=-\Delta c_n\geq0$；因此
$0\leq c_n\leq c_0=0$，零测度即为所求。以下设 $c_0>0$。

对每个 $N\geq1$ 及 $0\leq k\leq N$ 置
\begin{equation}\label{eq:7-3-hausdorff-weights}
 p_{N,k}=\binom Nk(-1)^{N-k}\Delta^{N-k}c_k,
 \qquad
 \mu_N=\sum_{k=0}^Np_{N,k}\delta_{k/N}.
\end{equation}
假设保证 $p_{N,k}\geq0$。为计算这些权重，不妨仍以 $L(x^j)=c_j$ 记系数泛函，并令
\[
 B_{N,k}(x)=\binom Nkx^k(1-x)^{N-k}
\]
为 Bernstein 基函数。展开 $(1-x)^{N-k}$ 可得
\begin{align*}
 L(B_{N,k})
 &=\binom Nk\sum_{j=0}^{N-k}(-1)^j\binom{N-k}{j}c_{k+j}\\
 &=\binom Nk(-1)^{N-k}\Delta^{N-k}c_k=p_{N,k}.
\end{align*}
Bernstein 恒等式 $\sum_{k=0}^NB_{N,k}=1$ 因而给
\[
 \sum_{k=0}^Np_{N,k}=L(1)=c_0.
\]
更精确地，记下降阶乘 $(k)_r=k(k-1)\cdots(k-r+1)$。二项分布的有限代数恒等式
\begin{equation}\label{eq:7-3-bernstein-factorial-identity}
 \sum_{k=0}^N(k)_rB_{N,k}(x)=(N)_rx^r
\end{equation}
可由
$(k)_r\binom Nk=(N)_r\binom{N-r}{k-r}$ 及二项式定理直接得到。对
式~\eqref{eq:7-3-bernstein-factorial-identity} 施加 $L$，便有
\begin{equation}\label{eq:7-3-hausdorff-factorial-moments}
 \sum_{k=0}^N(k)_rp_{N,k}=(N)_rc_r.
\end{equation}

令 $S(r,j)$ 为第二类 Stirling 数，即
$k^r=\sum_{j=0}^rS(r,j)(k)_j$。当 $N\geq r$ 时，
式~\eqref{eq:7-3-hausdorff-factorial-moments} 给出
\begin{align*}
 \int_0^1x^r\,\mu_N(\dd x)
 &=\frac1{N^r}\sum_{k=0}^Nk^rp_{N,k}\\
 &=\sum_{j=0}^rS(r,j)\frac{(N)_j}{N^r}c_j
 \longrightarrow c_r.
\end{align*}
最后一步使用 $S(r,r)=1$、$(N)_r/N^r\to1$，而 $j<r$ 时
$(N)_j/N^r\to0$。这一步说明普通幂矩不是被假定出来的，而是由下降阶乘矩逐阶恢复。

归一化测度 $c_0^{-1}\mu_N$ 全部支持在共同紧集 $[0,1]$ 上。再次应用
定理~\ref{theorem:5-3-3-2}，可取子列使
$c_0^{-1}\mu_{N_\ell}\Rightarrow\lambda$，其中 $\lambda$ 是概率测度。令 $\mu=c_0\lambda$，则
$\mu_{N_\ell}$ 对所有有界连续测试积分收敛于 $\mu$。每个 $x^r$ 在 $[0,1]$ 上有界连续，所以
\[
 \int_0^1x^r\,\mu(\dd x)
 =\lim_{\ell\to\infty}\int_0^1x^r\,\mu_{N_\ell}(\dd x)=c_r.
\]
存在性得证。

若 $\mu,\nu$ 是两个表示测度，则它们对每个多项式的积分相同。多项式由
Stone--Weierstrass 定理~\ref{theorem:1-3-3-2} 在 $C([0,1])$ 中一致稠密；又因两测度有限，
积分关于一致范数连续，故它们对每个连续函数的积分都相同。Radon 开集公式
\ref{theorem:3-1-1-4} 随即给出 $\mu=\nu$。唯一性得证。
\end{proof}

\begin{proposition}[GNS 原型]\label{proposition:7-3-1-4}
设 $L$ 是正矩泛函。公式
\[
 \langle p,q\rangle_L=L(\overline p\,q)
\]
在 $\C[x]$ 上定义半内积。其零空间
$\mathcal N=\{p:L(\overline p\,p)=0\}$ 在乘以 $x$ 后不变。因此
\[
 M_x[p]=[xp]
\]
在内积空间 $\C[x]/\mathcal N$ 上良定义；把该商空间完备化为 $\mathcal H_L$ 后，$M_x$ 是定义
在稠密子空间 $\C[x]/\mathcal N$ 上的对称算子。
\end{proposition}

\begin{proof}
因为矩 $c_n$ 为实数，$\langle p,q\rangle_L$ 是 Hermitian 半正定型。正性先给出
Cauchy--Schwarz 不等式。记
$a=L(\overline p p)$、$b=L(\overline q q)$、
$c=L(\overline p q)$。若 $b>0$，在
\[
 0\leq L(\overline{(p+\lambda q)}(p+\lambda q))
 =a+\lambda c+\overline\lambda\,\overline c+|\lambda|^2b
\]
中取 $\lambda=-\overline c/b$，便得
$0\leq a-|c|^2/b$，即
\[
 |L(\overline p q)|^2
 \leq L(\overline p p)L(\overline q q).
\]
若 $b=0$，则对 $t\in\R$ 分别取 $\lambda=t$ 与 $\lambda=it$，得到
\[
 0\leq a+2t\Re c,
 \qquad
 0\leq a-2t\Im c
 \qquad(t\in\R).
\]
让 $t$ 分别取任意大的正、负值，迫使 $\Re c=\Im c=0$。因此 $c=0$，同一不等式仍成立。

现在取 $p\in\mathcal N$。Cauchy--Schwarz 给
\[
 L(\overline p\,r)=0\qquad(r\in\C[x]).
\]
关键是不能只在这里声称零空间自动为理想；令 $r=x^2p$，才得到
\[
 L(\overline{xp}\,xp)=L(x^2\overline p\,p)
 =L(\overline p\,x^2p)=0.
\]
它还说明 $\mathcal N$ 是该半内积型的根空间，特别是线性子空间。上式给出
$xp\in\mathcal N$。若 $[p]=[q]$，则 $p-q\in\mathcal N$，从而
$x(p-q)\in\mathcal N$；所以 $M_x[p]=[xp]$ 与代表元无关。

最后，由 $x=\overline x$ 以及多项式乘法可交换，
\[
 \langle M_x[p],[q]\rangle_L
 =L(x\overline p\,q)
 =L(\overline p\,xq)
 =\langle[p],M_x[q]\rangle_L.
\]
这证明 $M_x$ 对称。商空间在其完备化中稠密，故定义域稠密。这里没有断言 $M_x$ 有界、闭或自伴；
这些性质以及自伴扩张所产生的谱测度属于后续泛函分析。
\end{proof}

\begin{figure}[htbp]
\centering
\begin{tikzpicture}[node distance=10mm and 18mm]
  \node[book strong node] (mom) {矩数据\\$(c_n)$};
  \node[book node,right=of mom] (mat) {有限截断\\$H_N$ 或 $T_N$};
  \node[book warm node,right=of mat] (pos) {二次型正性\\有限维检验};
  \node[book node,below=of mat] (ip) {半内积\\$\langle p,q\rangle_L$};
  \node[book node,below=of pos] (approx) {有限正测度\\$\mu_N$};
  \node[book strong node,below=of ip] (gns) {GNS 原型\\商空间与乘法算子};
  \node[book strong node,below=of approx] (meas) {表示测度\\$\mu$};
  \node[book warning node,below=of meas] (boundary) {支集与唯一性\\需另行检验};
  \draw[book arrow] (mom)--(mat);
  \draw[book arrow] (mat)--(pos);
  \draw[book arrow] (pos)--(ip);
  \draw[book arrow] (pos)--(approx);
  \draw[book dashed arrow] (ip)--(gns);
  \draw[book accent arrow] (approx)--node[right,book label]{紧性极限}(meas);
  \draw[book dashed arrow] (meas)--(boundary);
\end{tikzpicture}
\caption{矩问题的两条有限维构造途径。矩阵正性一方面产生 GNS 半内积，另一方面在 Herglotz 与 Hausdorff
问题中直接产生有限正测度；紧性只负责取极限，支集和唯一性仍需各自的检验。}
\label{fig:7-3-moment-positive-routes}
\end{figure}

这张图也解释了正定核与时间序列协方差的表示入口：每个有限样本只给出半正定矩阵；若这些矩阵彼此
相容，表示定理才把它们组织成一个测度或过程。正交多项式、\Pade{} 逼近与连分数则从 GNS 商空间中
对 $1,x,x^2,\ldots$ 做正交化开始。这些联系由上述有限维正性与紧性构造已经得到解释；
一般谱定理将在更完整的算子理论中给出它们的统一表述。

\begin{exercise}
设 $\mu=\sum_{\ell=1}^mw_\ell\delta_{a_\ell}$，其中权重为正、原子互异。证明
$\operatorname{rank}H_N=\min\{m,N+1\}$，并指出若允许权重有正有负，证明中的哪一步失效。
\end{exercise}
\begin{exercise}
验证 $[0,\infty)$ 支集给出移位 Hankel 正性，并构造一个 Hankel 半正定但移位 Hankel 不半正定的
矩序列。
\end{exercise}
\begin{exercise}
在 Herglotz 定理中直接计算 $P_N$ 的第 $n$ 个 Fourier 系数，并解释为何必须先把 $c_0=0$ 与
$c_0>0$ 分开。
\end{exercise}

矩表示回答“数据来自什么测度”。接下来转向另一个问题：给定一列已经存在的测度，哪些有限观测足以
保证它们有极限，并且极限没有在无穷远损失质量？

\subsection{Helly 选择、分布函数与弱收敛}\label{subsec:7-3-2}
\index{Helly 选择定理}\index{分布函数}\index{弱收敛}\index{模糊收敛}

实线上的概率测度可以压缩成一个递增函数。这个编码的优势是：在每个有理点取值都是一个有界数，
因此可以用可数对角法抽取子列；它的代价是跳点处不能要求普通点态收敛。移动原子把这一边界展示得
最清楚。

\begin{example}[移动 Dirac]\label{ex:7-3-2-1}
令 $\mu_n=\delta_{1/n}$。对每个 $f\in C_b(\R)$，
\[
 \int f\,\dd\mu_n=f(1/n)\longrightarrow f(0)=\int f\,\dd\delta_0,
\]
故 $\delta_{1/n}\Rightarrow\delta_0$。相应分布函数为
\[
 F_n(x)=\1_{[1/n,\infty)}(x),
 \qquad F(x)=\1_{[0,\infty)}(x).
\]
若 $x\neq0$，则对充分大的 $n$ 有 $F_n(x)=F(x)$；但
$F_n(0)=0$，而 $F(0)=1$。失败恰好发生在极限分布函数的跳点。
\end{example}

\begin{example}[质量逃逸]\label{ex:7-3-2-2}
令 $\mu_n=\delta_n$。若 $f\in C_c(\R)$ 且 $\supp f\subset[-R,R]$，则 $n>R$ 时
\[
 \int f\,\dd\mu_n=f(n)=0.
\]
所以 $\delta_n$ 模糊收敛到零测度。然而它不可能弱收敛到零测度，因为常数函数 $1$ 属于
$C_b(\R)$，而
\[
 \int1\,\dd\delta_n=1\not\longrightarrow0.
\]
事实上其分布函数在每个固定实数处都趋于零，候选极限的右端质量仍为零，不是概率分布函数。
\end{example}

\begin{example}[经验分布函数]\label{ex:7-3-2-3}
给定样本 $X_1,\ldots,X_n$，经验测度与经验分布函数分别是
\[
 \widehat\mu_n=\frac1n\sum_{k=1}^n\delta_{X_k},
 \qquad
 F_n(x)=\widehat\mu_n(( -\infty,x])
       =\frac1n\sum_{k=1}^n\1_{\{X_k\leq x\}}.
\]
若 $X_k$ 独立同分布且公共分布函数为 $F$，则对每个固定 $x$，随机变量
$\1_{\{X_k\leq x\}}$ 是均值 $F(x)$ 的 Bernoulli 变量，强大数律给
$F_n(x)\to F(x)$ a.s.。这里的“每个固定 $x$”不能直接换成“同时对所有 $x$”；后者是
Glivenko--Cantelli 定理所提供的更强一致结论。
\end{example}

\begin{definition}[分布函数]\label{def:7-3-2-1}
$\R$ 上有限正 Borel 测度 $\mu$ 的\emph{分布函数}是
\[
 F_\mu(x)=\mu(( -\infty,x]).
\]
\end{definition}

\begin{proposition}[分布函数的刻画]\label{proposition:7-3-2-distribution-functions}
有限正 Borel 测度的分布函数递增且右连续，并且
\[
 F_\mu(x-):=\lim_{y\uparrow x}F_\mu(y)=\mu(( -\infty,x)),
 \qquad
 F_\mu(x)-F_\mu(x-)=\mu(\{x\}).
\]
若 $\mu$ 是概率测度，则
\[
 \lim_{x\to-\infty}F_\mu(x)=0,
 \qquad
 \lim_{x\to+\infty}F_\mu(x)=1.
\]
反之，具有上述两个端点极限的右连续非减函数 $F$ 是唯一一个概率测度的分布函数。
\end{proposition}

\begin{proof}
单调性直接来自集合包含关系。若 $y_n\downarrow x$，则
$(-\infty,y_n]\downarrow(-\infty,x]$；有限测度从上连续，故 $F_\mu(y_n)\to F_\mu(x)$。
若 $y_n\uparrow x$，则
$(-\infty,y_n]\uparrow(-\infty,x)$；测度从下连续给出
$F_\mu(y_n)\to\mu((-\infty,x))$，继而由
$(-\infty,x]=(-\infty,x)\mathbin{\dot\cup}\{x\}$ 得到跳跃公式。令半直线
分别向空集与全空间单调收敛，便得概率情形的两个端点极限。

反过来，对这样的 $F$ 应用 Lebesgue--Stieltjes 构造
（定理~\ref{theorem:2-3-3-2}），得到唯一测度 $\mu_F$，满足
$\mu_F((a,b])=F(b)-F(a)$ 以及 $\mu_F(( -\infty,b])=F(b)$。再令 $b\to+\infty$ 可知
$\mu_F(\R)=1$，所以 $F$ 正是概率测度 $\mu_F$ 的分布函数。
\end{proof}

\begin{definition}[Helly 收敛]\label{def:7-3-2-2}
设 $F_n,F$ 都递增。若
\[
 F_n(x)\longrightarrow F(x)
\]
对 $F$ 的每个连续点 $x$ 成立，则称 $F_n$ 按连续点 \emph{Helly 收敛}于 $F$。
当这些函数编码有限测度时，还要另行核对 $\pm\infty$ 处的极限，以判断极限是否保留总质量。
\end{definition}

\begin{definition}[模糊收敛与弱收敛]\label{def:7-3-2-3}
设 $\mu_n,\mu$ 是 $\R$ 上局部有限正 Radon 测度。若
\[
 \int f\,\dd\mu_n\longrightarrow\int f\,\dd\mu
 \qquad(f\in C_c(\R)),
\]
则称 $\mu_n$ \emph{模糊收敛}（vague convergence）于 $\mu$。若 $\mu_n,\mu$ 是概率，并且同一收敛
对每个 $f\in C_b(\R)$ 成立，则称 $\mu_n$ \emph{弱收敛}（weak convergence）于 $\mu$，记作
$\mu_n\Rightarrow\mu$。从 $C_c$ 升到 $C_b$ 必须控制无穷远尾部；紧性正是这一统一控制。
\end{definition}

\begin{theorem}[Helly 选择定理]\label{theorem:7-3-2-1}
设 $(F_n)$ 是 $\R$ 上的递增函数序列，并存在有限常数 $A<B$ 使
$A\leq F_n(x)\leq B$ 对所有 $n,x$ 成立。则存在子列 $(F_{n_j})$ 与一个取值于 $[A,B]$ 的递增
右连续函数 $F$，使
\[
 F_{n_j}(x)\longrightarrow F(x)
\]
对 $F$ 的每个连续点成立。
\end{theorem}

\begin{proof}
枚举有理数 $\Q=\{q_1,q_2,\ldots\}$。有界数列 $(F_n(q_1))_n$ 有收敛子列；在该子列中再使
$F_n(q_2)$ 收敛，如此递推。取第 $j$ 层子列的第 $j$ 项所成对角子列，仍记为 $(F_{n_j})$，则
\[
 \phi(q)=\lim_{j\to\infty}F_{n_j}(q)
\]
对每个 $q\in\Q$ 存在。由于每个 $F_n$ 递增，$\phi$ 在 $\Q$ 上递增。

定义右连续规范
\begin{equation}\label{eq:7-3-helly-limit-definition}
 F(x)=\inf\{\phi(q):q\in\Q,\ q>x\}.
\end{equation}
$F$ 递增且仍取值于 $[A,B]$。若 $x_m\downarrow x$，则 $F(x_m)\geq F(x)$。给定
$\varepsilon>0$，由下确界定义可取有理数 $q>x$ 使
$\phi(q)<F(x)+\varepsilon$；当 $m$ 足够大时 $x_m<q$，从而
$F(x_m)\leq\phi(q)<F(x)+\varepsilon$。所以 $F(x_m)\downarrow F(x)$，即 $F$ 右连续。

对有理数 $a<x<b$，单调性给
\[
 F_{n_j}(a)\leq F_{n_j}(x)\leq F_{n_j}(b).
\]
令 $j\to\infty$，再分别对 $a\uparrow x$、$b\downarrow x$ 取上确界和下确界，得到
\begin{equation}\label{eq:7-3-helly-squeeze}
 F(x-)\leq\liminf_jF_{n_j}(x)
 \leq\limsup_jF_{n_j}(x)\leq F(x).
\end{equation}
这里
$F(x-)=\sup_{a<x,\,a\in\Q}\phi(a)$：一边由 $\phi$ 的单调性得到，另一边则在任意 $y<x$ 与
$x$ 之间插入一个有理数得到。若 $x$ 是 $F$ 的连续点，则 $F(x-)=F(x)$，
式~\eqref{eq:7-3-helly-squeeze} 迫使 $F_{n_j}(x)\to F(x)$。

最后说明这样的连续点足够多。对正整数 $m,r$，在 $[-m,m]$ 中跳跃大于 $1/r$ 的点只能有有限个：
若依次取有限多个这样的点，所有跳跃之和不超过 $B-A$，故点数至多 $r(B-A)+1$。任一不连续点都有
正跳跃，因而属于这些有限集的可数并。故递增函数的不连续点至多可数，证明完成。
\end{proof}

\begin{theorem}[分布函数判据]\label{theorem:7-3-2-2}
对 $\R$ 上概率测度 $\mu_n,\mu$，下列两件事等价：
\begin{enumerate}
  \item $\mu_n\Rightarrow\mu$；
  \item 对 $F_\mu$ 的每个连续点 $x$，有
  $F_{\mu_n}(x)\to F_\mu(x)$。
\end{enumerate}
\end{theorem}

\begin{proof}
先设 $\mu_n\Rightarrow\mu$。若 $x$ 是 $F_\mu$ 的连续点，则
$\mu(\{x\})=0$，所以集合 $(-\infty,x]$ 的边界对 $\mu$ 为零。Portmanteau 定理
\ref{theorem:5-3-3-1} 的连续集结论给
\[
 F_{\mu_n}(x)=\mu_n(( -\infty,x])\longrightarrow
 \mu(( -\infty,x])=F_\mu(x).
\]

反过来，设连续点收敛成立。先证明 $(\mu_n)$ 具有紧性。给定
$0<\varepsilon<1/4$，可选取
$F_\mu$ 的连续点 $a<b$，使
\[
 F_\mu(a)<\varepsilon,
 \qquad F_\mu(b)>1-\varepsilon;
\]
这是因为不连续点至多可数，而分布函数在两端趋于 $0,1$。对充分大的 $n$，
$F_{\mu_n}(a)<2\varepsilon$ 且 $F_{\mu_n}(b)>1-2\varepsilon$，故
\[
 \mu_n([a,b])\geq\mu_n((a,b])
 =F_{\mu_n}(b)-F_{\mu_n}(a)>1-4\varepsilon.
\]
有限多个尚未包括的 $\mu_n$ 各自可由一个紧区间捕获大于
$1-4\varepsilon$ 的质量；把这些区间与 $[a,b]$ 取有限并，便得到整列的统一紧区间。
由于 $4\varepsilon$ 可任意小，这正是 $(\mu_n)$ 的紧性。

 现取 $f\in C_b(\R)$。由刚证的紧性以及 $\mu$ 本身的紧性，可对任意
 $\eta>0$ 选 $R$，使
 \[
 \sup_n\mu_n([-R,R]^c)<\eta,
 \qquad \mu([-R,R]^c)<\eta.
 \]
 取 $\chi\in C_c(\R)$，满足 $0\leq\chi\leq1$ 且在 $[-R,R]$ 上等于一，并记
 $g=f\chi$。先直接由分布函数的区间增量证明
 $\int g\,\dd\mu_n\to\int g\,\dd\mu$。选取 $F_\mu$ 的连续点 $a'<b'$，使
 $\supp g\subset(a',b')$。给定
 $\delta>0$，利用 $g$ 在 $[a',b']$ 上的一致连续性，并利用 $F_\mu$ 的连续点在每个非空开区间中
 稠密，可取分割
 \[
  a'=x_0<x_1<\cdots<x_m=b'
 \]
 使所有 $x_j$ 都是 $F_\mu$ 的连续点，且 $g$ 在每个
 $[x_{j-1},x_j]$ 上的振幅小于 $\delta$。在每段选一点 $\xi_j$ 并令
 \[
  s=\sum_{j=1}^m g(\xi_j)\1_{(x_{j-1},x_j]}.
 \]
 因为 $\mu_n,\mu$ 都是概率测度，
 \[
  \left|\int(g-s)\,\dd\mu_n\right|\leq\delta,
  \qquad
  \left|\int(g-s)\,\dd\mu\right|\leq\delta.
 \]
 另一方面，
 \[
  \int s\,\dd\mu_n
  =\sum_{j=1}^m g(\xi_j)
    \bigl(F_{\mu_n}(x_j)-F_{\mu_n}(x_{j-1})\bigr)
  \longrightarrow\int s\,\dd\mu.
 \]
 因而
 \[
  \limsup_{n\to\infty}
  \left|\int g\,\dd\mu_n-\int g\,\dd\mu\right|\leq2\delta.
 \]
 令 $\delta\downarrow0$，得到所需的紧支测试收敛。另一方面，
 \begin{align*}
 \left|\int f(1-\chi)\,\dd\mu_n\right|
 &\leq\|f\|_\infty\mu_n([-R,R]^c),\\
 \left|\int f(1-\chi)\,\dd\mu\right|
 &\leq\|f\|_\infty\mu([-R,R]^c).
\end{align*}
先令 $n\to\infty$，再令 $\eta\downarrow0$，得到
$\int f\,\dd\mu_n\to\int f\,\dd\mu$。这正是弱收敛。
\end{proof}

\begin{proposition}[Stieltjes 积分收敛]\label{proposition:7-3-2-3}
设 $\mu_n,\mu$ 是 $\R$ 上有限正测度，分布函数分别为 $F_n,F$。假设
$F_n(x)\to F(x)$ 对 $F$ 的每个连续点成立。若 $f\in C_c(\R)$，并且在某个包含
$\supp f$ 的有界开区间 $(a,b)$ 上有
\[
 \sup_n\mu_n((a,b])<\infty,
\]
则
\[
 \int_\R f\,\dd F_n=\int_\R f\,\dd\mu_n
 \longrightarrow
 \int_\R f\,\dd\mu=\int_\R f\,\dd F.
\]
\end{proposition}

\begin{proof}
可把 $a,b$ 略微向区间内部移动，使它们都是 $F$ 的连续点，仍有
$\supp f\subset(a,b)$。记
$M=\sup_n\mu_n((a,b])$。给定 $\varepsilon>0$，由 $f$ 在 $[a,b]$ 上一致连续，可选取分割
\[
 a=x_0<x_1<\cdots<x_m=b,
\]
使每个 $x_j$ 都是 $F$ 的连续点，并且 $f$ 在每个 $[x_{j-1},x_j]$ 上的振幅小于
$\varepsilon$。这可以做到，因为 $F$ 的不连续点至多可数，连续点在每个区间中稠密。

在每段选 $\xi_j\in(x_{j-1},x_j)$，置
\[
 s(x)=\sum_{j=1}^mf(\xi_j)\1_{(x_{j-1},x_j]}(x).
\]
由于 $f$ 在 $(a,b]$ 外为零且 $f(a)=0$，
\begin{align*}
 \left|\int f\,\dd\mu_n-\int s\,\dd\mu_n\right|
 &\leq\varepsilon\mu_n((a,b])\leq\varepsilon M,\\
 \left|\int f\,\dd\mu-\int s\,\dd\mu\right|
 &\leq\varepsilon\mu((a,b]).
\end{align*}
而阶梯积分是有限和
\[
 \int s\,\dd\mu_n
 =\sum_{j=1}^mf(\xi_j)
   \bigl(F_n(x_j)-F_n(x_{j-1})\bigr),
\]
每个端点都是 $F$ 的连续点，所以此有限和趋于
$\int s\,\dd\mu$。于是
\[
 \limsup_{n\to\infty}
 \left|\int f\,\dd\mu_n-\int f\,\dd\mu\right|
 \leq\varepsilon\bigl(M+\mu((a,b])\bigr).
\]
令 $\varepsilon\downarrow0$ 即得结论。
\end{proof}

\begin{corollary}[有限正测度的分布函数判据]\label{corollary:7-3-2-finite-measures}
设 $\mu_n,\mu$ 是 $\R$ 上有限正 Borel 测度，且
$F_{\mu_n}(x)\to F_\mu(x)$ 对 $F_\mu$ 的每个连续点成立。则
\[
 \int f\dd\mu_n\longrightarrow\int f\dd\mu
 \qquad(f\in C_c(\R)).
\]
若还有 $\mu_n(\R)\to\mu(\R)$，则同一结论对每个 $f\in C_b(\R)$ 成立。
\end{corollary}

\begin{proof}
对 $f\in C_c(\R)$，可选 $F_\mu$ 的连续点 $a<b$ 使 $\supp f\subset(a,b)$。于是
\[
 \mu_n((a,b])=F_{\mu_n}(b)-F_{\mu_n}(a)
 \longrightarrow F_\mu(b)-F_\mu(a),
\]
故这些区间质量一致有界；命题~\ref{proposition:7-3-2-3} 给出
$\int f\dd\mu_n\to\int f\dd\mu$，不需要总质量收敛。

再设 $M_n=\mu_n(\R)\to M=\mu(\R)$。若 $M=0$，则 $\mu=0$，并且对每个
$f\in C_b(\R)$，
\[
 \left|\int f\dd\mu_n\right|\leq\norm f_\infty M_n\longrightarrow0.
\]
若 $M>0$，则 $M_n>0$ 最终成立。对这些 $n$ 令
$\bar\mu_n=M_n^{-1}\mu_n$、$\bar\mu=M^{-1}\mu$。它们是概率测度，而且在
$F_\mu$ 的每个连续点（也就是 $F_{\bar\mu}$ 的连续点）都有
\[
 F_{\bar\mu_n}(x)=\frac{F_{\mu_n}(x)}{M_n}\longrightarrow
 \frac{F_\mu(x)}M=F_{\bar\mu}(x).
\]
定理~\ref{theorem:7-3-2-2} 因而给出 $\bar\mu_n\Rightarrow\bar\mu$。对任意
$f\in C_b(\R)$，乘回总质量便得
\[
 \int f\dd\mu_n=M_n\int f\dd\bar\mu_n
 \longrightarrow M\int f\dd\bar\mu=\int f\dd\mu.
\]
这也给出了命题~\ref{proposition:2-3-3-4} 在有限正测度上的分布函数版本。
\end{proof}

\begin{corollary}[Helly--Prokhorov 对照]\label{corollary:7-3-2-4}
每个具有紧性的 $\R$ 上概率测度序列都有弱收敛子列。具体地，Helly 选择从其分布函数中提取连续点
收敛子列，紧性保证所得极限函数在 $-\infty,+\infty$ 的极限分别为 $0,1$。
\end{corollary}

\begin{proof}
设 $F_n$ 是 $\mu_n$ 的分布函数。它们递增且取值于 $[0,1]$，故定理
\ref{theorem:7-3-2-1} 给出子列 $F_{n_j}$ 与递增右连续函数 $F$，使子列在 $F$ 的连续点收敛。

给定 $\varepsilon>0$，紧性提供 $R>0$，使
$\mu_n([-R,R])\geq1-\varepsilon$ 对所有 $n$ 成立。若 $x<-R$ 是 $F$ 的连续点，则
\[
 F(x)=\lim_jF_{n_j}(x)\leq\varepsilon;
\]
若 $x>R$ 是连续点，则 $F(x)\geq1-\varepsilon$。连续点在实线两端都可任意选取，故
\[
 \lim_{x\to-\infty}F(x)\leq\varepsilon,
 \qquad
 \lim_{x\to+\infty}F(x)\geq1-\varepsilon.
\]
让 $\varepsilon\downarrow0$，再用 $0\leq F\leq1$，得到两端极限 $0,1$。因此由 $F$ 的
Lebesgue--Stieltjes 构造得到一个概率测度 $\mu$，其分布函数为 $F$。分布函数判据
\ref{theorem:7-3-2-2} 给 $\mu_{n_j}\Rightarrow\mu$。
\end{proof}

\begin{figure}[htbp]
\centering
\begin{tikzpicture}[x=1cm,y=1cm]
  \begin{scope}
    \draw[book line,-{Stealth[length=2mm]}] (-2.7,0)--(2.8,0) node[right]{$x$};
    \draw[book line,-{Stealth[length=2mm]}] (0,-.15)--(0,2.35) node[above]{$F$};
    \draw[BookAccent,very thick] (-2.5,0)--(0,0)--(0,2)--(2.5,2);
    \draw[BookWarm,semithick,dashed] (-2.5,0)--(.7,0)--(.7,2)--(2.5,2);
    \draw[BookWarning,semithick,dash pattern=on 1.2pt off 1.2pt]
      (-2.5,0)--(1.3,0)--(1.3,2)--(2.5,2);
    \fill[BookWarning] (0,0) circle (1.8pt);
    \fill[BookAccent] (0,2) circle (1.8pt);
    \node[book label,align=center] at (0,-.65) {跳点 $x=0$：\\不要求点态收敛};
    \node[book label] at (-1.45,2.5) {移动 Dirac};
  \end{scope}
  \begin{scope}[xshift=8cm]
    \draw[book line,-{Stealth[length=2mm]}] (-2.8,0)--(3,0) node[right]{$x$};
    \fill[BookAccent!10] (-1.4,-.35) rectangle (1.4,.35);
    \draw[BookAccent,semithick] (-1.4,-.35) rectangle (1.4,.35);
    \node[book label] at (0,-.72) {固定紧区间};
    \foreach \x/\lab in {1.8/$\delta_n$,2.25/$\delta_{n+1}$,2.7/$\delta_{n+2}$}{
      \fill[BookWarning] (\x,0) circle (2.2pt);
      \node[book label,rotate=55] at (\x,.55) {\lab};
    }
    \draw[book dashed arrow] (1.65,1.15)--(2.8,1.15)
      node[midway,above,book label]{质量逃向 $+\infty$};
    \node[book label] at (0,2.15) {$C_c$ 看不见图外质量};
  \end{scope}
\end{tikzpicture}
\caption{左：移动原子的分布函数只在极限跳点失配。右：$\delta_n$ 离开每个固定紧区间，因而可
模糊收敛到零，却不能作为概率弱收敛。}
\label{fig:7-3-helly-jump-escape}
\end{figure}

Helly 选择是一般 Prokhorov 定理在实线上的坐标化版本：有理点对角抽取负责选出候选，
紧性负责保证候选仍有全部质量，分布函数判据负责把坐标收敛还原为弱收敛。经验分布、
统计量的分布与 Stieltjes 变换的极限都可沿这条路线处理。一般 Polish 空间没有一条天然的实数坐标，
因而要使用第~5~章的可数有界 Lipschitz 测试族；在更一般的拓扑空间中，序列甚至可能不够，必须回到
第~1~章的网。

\begin{exercise}
证明递增函数在有界区间内跳跃大于 $\varepsilon$ 的点只有有限个，并由此重新证明其不连续点至多
可数。
\end{exercise}
\begin{exercise}
逐点计算 $\delta_{1/n}$ 的分布函数，并分别在 $x<0$、$x=0$、$x>0$ 检验
定理~\ref{theorem:7-3-2-2}。
\end{exercise}
\begin{exercise}
设有限正测度 $\mu_n$ 模糊收敛于 $\mu$，且 $\mu_n(\R)\to\mu(\R)$。用紧支撑截断证明
$\int f\,\dd\mu_n\to\int f\,\dd\mu$ 对每个 $f\in C_b(\R)$ 成立。
\end{exercise}

分布函数以半直线指标观测测度。下面改用复指数、指数衰减核、Cauchy 核与多项式，并逐一说明哪些
测试族总能识别测度，哪些测试族需要额外增长条件。

\subsection{变换决定测度、矩确定性与 Carleman 条件}\label{subsec:7-3-3}
\index{Carleman 条件}\index{矩方法}\index{Laplace 变换}\index{准解析性}

特征函数对每个概率都存在，并且总能决定分布；Laplace 与 Cauchy--Stieltjes 变换只在各自的自然域
定义，却把支集或预解式结构直接编码进去；多项式矩最便于代数计算，但在无界支集上可能无法
识别测度。差别不在于这些变换是否“形式上含有全部系数”，而在于局部数据能否通过增长控制传播到
整个定义域。

\begin{example}[紧支集确定性]\label{ex:7-3-3-1}
设有限正测度 $\mu,\nu$ 都支持在共同紧区间 $K=[-R,R]$，并满足
\[
 \int_Kx^n\,\mu(\dd x)=\int_Kx^n\,\nu(\dd x)
 \qquad(n=0,1,2,\ldots).
\]
线性组合说明两测度对每个多项式积分相同。给定 $f\in C(K)$，由 Stone--Weierstrass 定理可取
多项式 $p_j$ 使 $\|p_j-f\|_{\infty,K}\to0$，于是
\begin{align*}
 \left|\int_K f\,\dd(\mu-\nu)\right|
 &\leq\left|\int_K(f-p_j)\,\dd\mu\right|
      +\left|\int_K(f-p_j)\,\dd\nu\right|\\
 &\leq\|f-p_j\|_{\infty,K}\bigl(\mu(K)+\nu(K)\bigr)\longrightarrow0.
\end{align*}
故两测度对全部连续测试相同，从而 $\mu=\nu$。紧支集同时保证多项式有界并允许一致逼近；这两点在
$\R$ 上都不再自动成立。
\end{example}

\begin{example}[Gaussian Carleman]\label{ex:7-3-3-2}
标准 Gaussian 的奇矩为零，偶矩满足
\[
 m_{2n}=\E X^{2n}=(2n-1)!!=\frac{(2n)!}{2^nn!}.
\]
Stirling 公式给
\[
 m_{2n}^{1/(2n)}\sim\sqrt{\frac{2n}{e}},
 \qquad
 m_{2n}^{-1/(2n)}\sim\sqrt{\frac{e}{2n}}.
\]
因此 $\sum_nm_{2n}^{-1/(2n)}$ 与 $\sum_nn^{-1/2}$ 一样发散。即使不用 Stirling，也有
$m_{2n}=\prod_{j=1}^n(2j-1)\leq(2n)^n$，从而
$m_{2n}^{-1/(2n)}\geq(2n)^{-1/2}$。下文的 Carleman 定理于是说明 Gaussian 分布由全部矩唯一
决定。
\end{example}

\begin{example}[矩收敛但质量问题]\label{ex:7-3-3-3}
固定整数 $r\geq1$，对 $R>1$ 令
\[
 \mu_R=(1-R^{-(r+1)})\delta_0+R^{-(r+1)}\delta_R.
\]
其总质量为一。对 $1\leq k\leq r$，
\[
 \int x^k\,\dd\mu_R=R^{k-r-1}\longrightarrow0
 =\int x^k\,\dd\delta_0,
\]
而
\[
 \int x^{r+1}\,\dd\mu_R=1.
\]
事实上 $\mu_R\Rightarrow\delta_0$，因为对 $f\in C_b(\R)$，
\[
 \left|\int f\,\dd\mu_R-f(0)\right|
 =R^{-(r+1)}|f(R)-f(0)|
 \leq2\|f\|_\infty R^{-(r+1)}\to0.
\]
所以这个例子揭示的是“任意固定的前 $r$ 阶矩不能控制下一阶矩”，不是紧性失败。若改取
$\delta_R$，远端质量不衰减，才会连紧性也失败。两种尾部机制必须分开。
\end{example}

\begin{definition}[Laplace 变换]\label{def:7-3-3-1}
对 $[0,\infty)$ 上有限正 Borel 测度 $\mu$，定义
\[
 L_\mu(s)=\int_{[0,\infty)}e^{-sx}\,\mu(\dd x),\qquad s\geq0.
\]
同一公式在复半平面 $\Re z>0$ 给出全纯函数。若
$\int e^{a x}\,\dd\mu<\infty$ 对某个 $a>0$ 成立，则定义域可越过零延伸到
$\Re z>-a$，并可在零附近逐阶求导。
\end{definition}

\begin{definition}[Carleman 条件]\label{def:7-3-3-2}
设概率测度 $\mu$ 的全部偶矩
$m_{2n}=\int_\R x^{2n}\,\mu(\dd x)$ 有限。若
\begin{equation}\label{eq:7-3-carleman-condition}
 \sum_{n=1}^\infty m_{2n}^{-1/(2n)}=\infty,
\end{equation}
则称 $\mu$ 满足 \emph{Carleman 条件}。若某个 $m_{2n}=0$，相应项按 $+\infty$ 解释；此时
$\mu=\delta_0$。
\end{definition}

\begin{definition}[矩方法的标准假设]\label{def:7-3-3-3}
设 $X_n,X$ 为实随机变量。称它们满足\emph{矩方法的标准假设}，若：
\begin{enumerate}
  \item 对每个整数 $k\geq0$，$\E|X_n|^k$ 与 $\E|X|^k$ 都有限；
  \item 对每个 $k\geq0$，有 $\E X_n^k\to m_k=\E X^k$；
  \item 矩序列 $(m_k)$ 在 $\R$ 上确定 $X$ 的概率律。
\end{enumerate}
这里必须给出所有阶矩，而不是一个固定的有限截断。证明紧性使用二阶矩，传递第 $k$ 阶矩则
使用一个严格更高的偶数阶矩。
\end{definition}

\begin{theorem}[变换唯一性链]\label{theorem:7-3-3-1}
下列三类变换都决定相应测度。
\begin{enumerate}
  \item $\R$ 上两个有限复 Borel 测度若具有相同的 Fourier--Stieltjes 变换
  $\widehat\mu(t)=\int e^{itx}\,\mu(\dd x)$，则两测度相同。
  \item $\R$ 上两个有限实符号测度若其 Cauchy 变换
  $G_\mu(z)=\int(z-x)^{-1}\,\mu(\dd x)$ 在上半平面相同，则两测度相同。对一般有限复测度，给定
  上、下两个半平面的相同数据仍有同一结论。
  \item $[0,\infty)$ 上两个有限实符号测度若其 Laplace 变换对所有 $s\geq0$ 相同，则两测度
  相同。
\end{enumerate}
\end{theorem}

\begin{proof}
每一项都对两测度之差 $\sigma$ 证明：相应变换为零蕴含 $\sigma=0$。

先证 Fourier 情形。令
\[
 \gamma_\varepsilon(x)=\frac1{\sqrt{4\pi\varepsilon}}
 \exp\!\left(-\frac{x^2}{4\varepsilon}\right),
 \qquad \varepsilon>0.
\]
Gaussian 反演公式是
\[
 \gamma_\varepsilon(u)=\frac1{2\pi}
 \int_\R e^{itu}e^{-\varepsilon t^2}\,\dd t.
\]
由于 $|\sigma|(\R)<\infty$ 且 $e^{-\varepsilon t^2}\in L^1(\R)$，Fubini 定理给卷积测度一个连续
密度
\begin{align*}
 (\gamma_\varepsilon*\sigma)(x)
 &=\int_\R\gamma_\varepsilon(x-y)\,\sigma(\dd y)\\
 &=\frac1{2\pi}\int_\R e^{itx}e^{-\varepsilon t^2}
   \widehat\sigma(-t)\,\dd t=0.
\end{align*}
对 $f\in C_c(\R)$ 再用 Fubini，
\[
 \int_\R(f*\gamma_\varepsilon)(y)\,\sigma(\dd y)
 =\int_\R f(x)(\gamma_\varepsilon*\sigma)(x)\,\dd x=0.
\]
Gaussian 是近似恒等核，所以 $f*\gamma_\varepsilon\to f$ 一致；乘以有限全变差并令
$\varepsilon\downarrow0$，得到 $\int f\,\dd\sigma=0$。有限复 Radon 测度由 $C_c$ 测试唯一决定，
故 $\sigma=0$。

再证 Cauchy 情形。记 Poisson 核
\[
 P_v(u)=\frac1\pi\frac{v}{u^2+v^2},\qquad v>0.
\]
若 $\sigma$ 为实符号测度，则
\begin{equation}\label{eq:7-3-cauchy-poisson-boundary}
 -\frac1\pi\Im G_\sigma(u+iv)
 =\int_\R P_v(u-x)\,\sigma(\dd x)=(P_v*\sigma)(u).
\end{equation}
所以上半平面中 $G_\sigma=0$ 蕴含 $P_v*\sigma=0$ 对每个 $v>0$ 成立。Poisson 核总积分为一，
并且对每个 $\delta>0$，
$\int_{|u|>\delta}P_v(u)\,\dd u\to0$。把积分分成 $|u|<\delta$ 与其补集，利用
$f\in C_c\subset C_0$ 的一致连续性，可得 $f*P_v\to f$ 一致。因此
\[
 \int f\,\dd\sigma
 =\lim_{v\downarrow0}\int(f*P_v)\,\dd\sigma
 =\lim_{v\downarrow0}\int f(u)(P_v*\sigma)(u)\,\dd u=0.
\]
故 $\sigma=0$。若 $\sigma$ 为复测度，则不能把虚部移入复测度积分；但两个半平面的数据给出
\[
 G_\sigma(u+iv)-G_\sigma(u-iv)
 =-2\pi i(P_v*\sigma)(u).
\]
若两边的 Cauchy 数据都为零，仍有 $P_v*\sigma=0$，上述 $C_c$ 逼近对复测度同样成立。

最后证 Laplace 情形。令 $T:[0,\infty)\to(0,1]$，$T(x)=e^{-x}$，并把推前符号测度
$\tau=T_*\sigma$ 视为 $[0,1]$ 上在点 $0$ 无质量的测度。对每个整数 $n\geq0$，
\[
 \int_{[0,1]}y^n\,\tau(\dd y)
 =\int_{[0,\infty)}e^{-nx}\,\sigma(\dd x)=L_\sigma(n)=0.
\]
故 $\tau$ 对所有多项式积分为零。多项式在 $C([0,1])$ 中一致稠密，有限全变差使积分关于一致范数
连续，所以 $\tau$ 对所有连续函数积分为零，进而 $\tau=0$。映射 $T$ 是
$[0,\infty)$ 与 $(0,1]$ 间的 Borel 双射，故 $\sigma=0$。三项均证毕。
\end{proof}

Carleman 条件比指数矩条件弱，不能把 Fourier 变换当成零点附近的全纯函数。所需传播原理是一个
纯实变量的准解析引理。下面只证明本应用需要的对数凸特例；矩产生的序列本身就对数凸，不需要把任意
序列替换成某个不保导数上界的“正则化”。

\begin{lemma}[Carleman--Bang 实变量准解析引理]\label{lemma:7-3-3-carleman-bang}
设 $M_0=1$、$M_n>0$，并且
\[
 M_n^2\leq M_{n-1}M_{n+1}\qquad(n\geq1).
\]
设 $F\in C^\infty(\R;\C)$，且存在常数 $C,A>0$ 使
\[
 \|F^{(n)}\|_\infty\leq CA^nM_n\qquad(n\geq0).
\]
若
$\sum_{n\geq1}M_n^{-1/n}=\infty$，并且在某个 $t_0\in\R$ 有
$F^{(n)}(t_0)=0$ 对所有 $n$ 成立，则 $F\equiv0$。
\end{lemma}

\begin{proof}
平移自变量，并以
\[
 \widetilde F(s)=C^{-1}F(t_0+s/A)
\]
替换 $F$，可把 $t_0=0$ 以及 $C=A=1$ 作为假设。以下省略波浪号。令
\[
 \mu_n=\frac{M_n}{M_{n-1}},\qquad n\geq1.
\]
对数凸性等价于 $\mu_n$ 递增。

第一步由题设级数推出所需的比值级数发散。置 $a_n=\mu_n^{-1}$，则因 $M_0=1$，
\[
 M_n^{-1/n}=(a_1a_2\cdots a_n)^{1/n}=:G_n.
\]
对 $a_1,2a_2,\ldots,na_n$ 用算术--几何平均不等式，得到
\[
 G_n\leq\frac{\sum_{j=1}^nj a_j}{n(n!)^{1/n}}.
\]
又由
$\log(n!)=\sum_{j=1}^n\log j\geq\int_1^n\log x\,\dd x\geq n\log n-n$
可得 $(n!)^{1/n}\geq n/e$，所以
\begin{equation}\label{eq:7-3-carleman-amgm}
 G_n\leq\frac{e}{n^2}\sum_{j=1}^nj a_j.
\end{equation}
对任意 $N$ 把有限和换序，并使用
$j\sum_{n=j}^\infty n^{-2}\leq2$，便有
\[
 \sum_{n=1}^NG_n
 \leq e\sum_{j=1}^Nj a_j\sum_{n=j}^N\frac1{n^2}
 \leq2e\sum_{j=1}^Na_j.
\]
因此题设中 $\sum G_n=\infty$ 蕴含
\begin{equation}\label{eq:7-3-carleman-ratio-divergence}
 \sum_{n=1}^\infty\frac1{\mu_n}=\infty.
\end{equation}

第二步定义 Bang 层级函数
\begin{equation}\label{eq:7-3-bang-function}
 B(x)=\sup_{n\geq0}\frac{|F^{(n)}(x)|}{e^nM_n}.
\end{equation}
归一化导数界给 $0\leq B\leq1$。第 $n$ 项统一不超过 $e^{-n}$，故
$B$ 是有限个连续函数之最大值的统一极限：只取 $0\leq n<N$ 所产生的最大值与 $B$ 相差至多
$e^{-N}$。因此 $B$ 连续。所有导数在零点消失又给 $B(0)=0$。

第三步建立传播估计。固定整数 $q\geq1$ 以及 $x,x+h\in\R$。若 $0\leq n<q$，对
$F^{(n)}$ 写到 $q-n-1$ 阶并使用积分余项：
\begin{align*}
 F^{(n)}(x+h)
 &=\sum_{\ell=0}^{q-n-1}
   \frac{F^{(n+\ell)}(x)}{\ell!}h^\ell\\
 &\quad+\frac{h^{q-n}}{(q-n-1)!}
   \int_0^1(1-s)^{q-n-1}F^{(q)}(x+sh)\,\dd s.
\end{align*}
由于 $\mu_j$ 递增，$n+\ell\leq q$ 时
\[
 \frac{M_{n+\ell}}{M_n}\leq\mu_q^\ell,
 \qquad
 \frac{M_q}{M_n}\leq\mu_q^{q-n}.
\]
结合式~\eqref{eq:7-3-bang-function} 与
$\|F^{(q)}\|_\infty\leq M_q$，得到
\begin{align*}
 \frac{|F^{(n)}(x+h)|}{e^nM_n}
 &\leq B(x)\sum_{\ell=0}^{q-n-1}
   \frac{(e\mu_q|h|)^\ell}{\ell!}\\
 &\quad+e^{-q}\frac{(e\mu_q|h|)^{q-n}}{(q-n)!}\\
 &\leq\max\{B(x),e^{-q}\}\exp(e\mu_q|h|).
\end{align*}
若 $n\geq q$，左端不超过 $e^{-n}\leq e^{-q}$。对 $n$ 取上确界，得到 Bang 基本不等式
\begin{equation}\label{eq:7-3-bang-propagation}
 B(x+h)\leq\max\{B(x),e^{-q}\}\exp(e\mu_q|h|).
\end{equation}

最后用层级穿越反证。若某个 $t>0$ 满足 $B(t)>0$，可选整数 $K$ 使
$e^{-K}<B(t)$。对每个 $k\geq K$ 定义
\[
 x_k=\inf\{x\in[0,t]:B(x)\geq e^{-k}\}.
\]
集合非空，因为 $B(t)>e^{-K}\geq e^{-k}$。又因 $B(0)=0<e^{-k}$ 及 $B$ 连续，有
$B(x_k)=e^{-k}$。层级 $e^{-(k+1)}$ 小于 $e^{-k}$，所以
$0\leq x_{k+1}\leq x_k\leq t$。把
式~\eqref{eq:7-3-bang-propagation} 用于
$x=x_{k+1}$、$h=x_k-x_{k+1}$、$q=k+1$，得到
\[
 e^{-k}\leq e^{-(k+1)}
 \exp\!\bigl(e\mu_{k+1}(x_k-x_{k+1})\bigr),
\]
因而
\[
 x_k-x_{k+1}\geq\frac1{e\mu_{k+1}}.
\]
对 $k=K,\ldots,N$ 求和可得
\[
 t\geq x_K-x_{N+1}
 \geq\frac1e\sum_{k=K}^N\frac1{\mu_{k+1}},
\]
这与式~\eqref{eq:7-3-carleman-ratio-divergence} 矛盾。故 $B(t)=0$ 对所有 $t>0$。对函数
$s\mapsto F(-s)$ 重复同一论证，得到负半轴上的结论。于是 $B\equiv0$，而
$|F(x)|/M_0\leq B(x)$，所以 $F\equiv0$。整个论证只用了实变量 Taylor 公式、AM--GM 与连续
层级穿越，没有使用复分析恒等定理。
\end{proof}

\begin{theorem}[Carleman 定理]\label{theorem:7-3-3-2}
若 $\R$ 上概率测度 $\mu$ 满足 Carleman 条件
\eqref{eq:7-3-carleman-condition}，则 $\mu$ 的 Hamburger 矩问题确定：任何与 $\mu$ 具有全部相同
矩的概率测度都等于 $\mu$。
\end{theorem}

\begin{proof}
先证明一个核心稠密性结论。设 $\rho$ 是满足 Carleman 条件的概率测度。若某个
$m_{2n}=\int|x|^{2n}\,\dd\rho$ 为零，则 $\rho$ 集中在 $\{0\}$，即 $\rho=\delta_0$，多项式在
$L^2(\rho)$ 中当然稠密。以下设所有偶矩为正，并置
\[
 M_n=m_{2n}^{1/2},\qquad M_0=1.
\]
Cauchy--Schwarz 不等式给
\[
 m_{2n}^2
 =\left(\int |x|^{n-1}|x|^{n+1}\,\rho(\dd x)\right)^2
 \leq m_{2n-2}m_{2n+2}.
\]
因此
\[
 M_n^2=m_{2n}
 \leq\sqrt{m_{2n-2}m_{2n+2}}=M_{n-1}M_{n+1};
\]
也就是说，所需的序列自身对数凸。Carleman 级数恰为
$\sum M_n^{-1/n}=\infty$。

设 $f\in L^2(\rho)$ 与所有多项式正交。若 $\norm f_{L^2(\rho)}=0$，则 $f=0$ 已经成立；以下设该
范数为正，并定义
\[
 F(t)=\int_\R e^{itx}f(x)\,\rho(\dd x).
\]
对每个 $n$，Cauchy--Schwarz 给
\[
 \int|x|^n|f(x)|\,\rho(\dd x)
 \leq\|f\|_{L^2(\rho)}m_{2n}^{1/2}<\infty.
\]
由 $|(e^{ihx}-1)/h|\leq|x|$ 反复应用支配收敛，可逐阶求导：
\[
 F^{(n)}(t)=i^n\int_\R x^ne^{itx}f(x)\,\rho(\dd x).
\]
于是
\[
 \|F^{(n)}\|_\infty
 \leq\|f\|_{L^2(\rho)}M_n,
 \qquad
 F^{(n)}(0)=i^n\int x^nf\,\dd\rho=0.
\]
引理~\ref{lemma:7-3-3-carleman-bang} 给出 $F\equiv0$。复测度
$\sigma=f\rho$ 有限，因为
$\int|f|\,\dd\rho\leq\|f\|_2$；其 Fourier--Stieltjes 变换正是 $F$。由
定理~\ref{theorem:7-3-3-1} 的第一项，$\sigma=0$，故 $f=0$ $\rho$-a.e.。这证明多项式在
$L^2(\rho)$ 中稠密。

现在设概率 $\nu$ 与 $\mu$ 的全部矩相同，并令 $\rho=(\mu+\nu)/2$。$\rho$ 与 $\mu$ 有同样的
偶矩，所以也满足 Carleman 条件。两测度都绝对连续于 $\rho$；令
\[
 h=\frac{\dd(\mu-\nu)}{\dd\rho}.
\]
由 $|\mu-\nu|\leq\mu+\nu=2\rho$ 可知 $|h|\leq2$ $\rho$-a.e.，故 $h\in L^2(\rho)$。对每个
$n\geq0$，
\[
 \int x^nh\,\dd\rho=\int x^n\,\dd\mu-\int x^n\,\dd\nu=0.
\]
刚证的多项式稠密性迫使 $h=0$，所以 $\mu-\nu=0$。这就证明了矩确定性。
\end{proof}

例~\ref{ex:7-3-3-2} 验证 Gaussian 满足条件；例~\ref{ex:7-3-1-3} 的对数正态 Stieltjes 族则
展示了条件缺失时的真实失败。准解析引理并未预设 $F$ 具有全纯延拓；它仅凭导数上界与一个
发散级数，把一点的无限阶平坦性传播到整条实线。

\begin{proposition}[指数矩蕴含确定性]\label{proposition:7-3-3-3}
若 $\R$ 上概率测度 $\mu$ 满足
\[
 \int_\R e^{a|x|}\,\mu(\dd x)<\infty
\]
对某个 $a>0$ 成立，则 $\mu$ 由全部矩唯一决定。
\end{proposition}

\begin{proof}
设概率 $\nu$ 与 $\mu$ 具有相同矩。由非负级数的单调收敛，
\begin{align*}
 \int_\R\cosh(ax)\,\nu(\dd x)
 &=\sum_{n=0}^\infty\frac{a^{2n}}{(2n)!}
   \int x^{2n}\,\nu(\dd x)\\
 &=\sum_{n=0}^\infty\frac{a^{2n}}{(2n)!}
   \int x^{2n}\,\mu(\dd x)
 =\int_\R\cosh(ax)\,\mu(\dd x)<\infty.
\end{align*}
由于 $e^{a|x|}\leq2\cosh(ax)$，$\nu$ 也有 $a$ 阶双边指数矩。

因此
\[
 M_\mu(z)=\int e^{zx}\,\mu(\dd x),
 \qquad
 M_\nu(z)=\int e^{zx}\,\nu(\dd x)
\]
都在带形域 $|\Re z|<a$ 全纯。对 $|z|<a$，$e^{zx}$ 的幂级数由 $e^{|z||x|}$ 支配，可以逐项
积分，故
\[
 M_\mu(z)=\sum_{n=0}^\infty\frac{z^n}{n!}\int x^n\,\dd\mu
 =\sum_{n=0}^\infty\frac{z^n}{n!}\int x^n\,\dd\nu=M_\nu(z).
\]
全纯恒等定理把相等性扩展到整个带形域；特别地，在 $z=it$ 上两者的特征函数相同。
定理~\ref{theorem:7-3-3-1} 的第一项给 $\mu=\nu$。这里全纯性由指数矩真正保证；Carleman 定理
没有这一保证，因而不能用这段证明代替准解析论证。这与评注
\ref{rem:5-4-1-complex-mgf} 所强调的复参数矩母函数边界一致。
\end{proof}

\begin{theorem}[矩方法的收敛定理]\label{theorem:7-3-3-4}
若实随机变量 $X_n,X$ 满足定义~\ref{def:7-3-3-3} 中的矩方法标准假设，则
\[
 X_n\Rightarrow X.
\]
\end{theorem}

\begin{proof}
记 $\mu_n=\Law(X_n)$、$\mu=\Law(X)$。二阶矩收敛给出
\[
 C_2:=\sup_n\int x^2\,\mu_n(\dd x)<\infty.
\]
Markov 不等式于是给
\[
 \sup_n\mu_n([-R,R]^c)
 \leq\frac{C_2}{R^2}\longrightarrow0.
\]
所以测度族 $(\mu_n)$ 具有紧性。

任取一个子列。由 Prokhorov 定理~\ref{theorem:5-3-3-2}，可再取子列
$(\mu_{n_j})$ 弱收敛于某个概率测度 $\nu$。下面逐阶证明 $\nu$ 保留全部矩。固定 $k\geq0$，取偶数
$q>k$。$q$ 阶矩收敛且 $q$ 为偶数，故
\[
 C_q:=\sup_n\int|x|^q\,\mu_n(\dd x)<\infty.
\]
对 $R>0$，
\begin{equation}\label{eq:7-3-method-ui-original}
 \sup_n\int_{|x|>R}|x|^k\,\mu_n(\dd x)
 \leq R^{k-q}C_q\longrightarrow0.
\end{equation}
这正是第 $k$ 阶矩所需的一致可积尾估计。

还必须先给弱极限建立同样的高阶界。对
$g_M(x)=\min\{|x|^q,M\}$，函数 $g_M$ 有界连续，所以
\[
 \int g_M\,\dd\nu
 =\lim_{j\to\infty}\int g_M\,\dd\mu_{n_j}\leq C_q.
\]
令 $M\uparrow\infty$ 并用单调收敛，得到
$\int|x|^q\,\dd\nu\leq C_q$，从而
\begin{equation}\label{eq:7-3-method-ui-limit}
 \int_{|x|>R}|x|^k\,\nu(\dd x)
 \leq R^{k-q}C_q\longrightarrow0.
\end{equation}

取连续截断 $0\leq\chi_R\leq1$，使 $\chi_R=1$ 于 $[-R,R]$，并在某个更大的紧区间外为零。
函数 $x^k\chi_R(x)$ 有界连续，故
\[
 \int x^k\chi_R(x)\,\mu_{n_j}(\dd x)
 \longrightarrow
 \int x^k\chi_R(x)\,\nu(\dd x).
\]
式~\eqref{eq:7-3-method-ui-original} 与
式~\eqref{eq:7-3-method-ui-limit} 控制等式两边去掉截断的误差。先令 $j\to\infty$，再令
$R\to\infty$，便得到
\[
 \int x^k\,\nu(\dd x)
 =\lim_{j\to\infty}\int x^k\,\mu_{n_j}(\dd x)
 =m_k=\int x^k\,\mu(\dd x).
\]
由于 $k$ 任意，$\nu$ 与 $\mu$ 的全部矩相同；标准假设中的矩确定性给出 $\nu=\mu$。

我们已经证明：原序列的每个子列都有进一步子列弱收敛到 $\mu$。若 $\mu_n$ 本身不弱收敛到
$\mu$，则由弱拓扑的有界 Lipschitz 度量化命题~\ref{proposition:5-3-3-weak-metric}，存在
$\varepsilon>0$ 与子列使
$d_{\mathrm{BL}}(\mu_{n_j},\mu)\geq\varepsilon$；但该子列又有进一步子列弱收敛于 $\mu$，与
$d_{\mathrm{BL}}\to0$ 矛盾。因此 $\mu_n\Rightarrow\mu$，即 $X_n\Rightarrow X$。
\end{proof}

\begin{figure}[htbp]
\centering
\begin{tikzpicture}[node distance=5.5mm and 15mm]
  \node[book strong node,minimum width=3.2cm] (h1) {测试族};
  \node[book strong node,right=of h1,minimum width=3.7cm] (h2) {自然定义域／代价};
  \node[book strong node,right=of h2,minimum width=3.7cm] (h3) {识别结论};

  \node[book node,below=of h1,minimum width=3.2cm] (cb) {$C_b$};
  \node[book node,right=of cb,minimum width=3.7cm] (cbd) {概率上总可测试};
  \node[book warm node,right=of cbd,minimum width=3.7cm] (cbr) {定义弱收敛并决定测度};

  \node[book node,below=of cb,minimum width=3.2cm] (cf) {$e^{itx}$};
  \node[book node,right=of cf,minimum width=3.7cm] (cfd) {对有限测度总存在};
  \node[book warm node,right=of cfd,minimum width=3.7cm] (cfr) {无条件唯一};

  \node[book node,below=of cf,minimum width=3.2cm] (cs) {$(z-x)^{-1}$};
  \node[book node,right=of cs,minimum width=3.7cm] (csd) {$z$ 离开实支集};
  \node[book warm node,right=of csd,minimum width=3.7cm] (csr) {Poisson 边界唯一};

  \node[book node,below=of cs,minimum width=3.2cm] (la) {$e^{-sx}$};
  \node[book node,right=of la,minimum width=3.7cm] (lad) {$[0,\infty)$，$s\geq0$};
  \node[book warm node,right=of lad,minimum width=3.7cm] (lar) {Stone--Weierstrass 唯一};

  \node[book node,below=of la,minimum width=3.2cm] (mo) {$1,x,x^2,\ldots$};
  \node[book warning node,right=of mo,minimum width=3.7cm] (mod) {矩可能不存在或增长过快};
  \node[book warning node,right=of mod,minimum width=3.7cm] (mor) {紧支集／Carleman 时唯一};

  \foreach \a/\b in {cb/cbd,cbd/cbr,cf/cfd,cfd/cfr,cs/csd,csd/csr,la/lad,lad/lar,mo/mod}{
    \draw[book arrow] (\a)--(\b);
  }
  \draw[book dashed arrow,draw=BookWarning] (mod)--node[above=1mm,book label,font=\scriptsize]{附加条件}(mor);
\end{tikzpicture}
\caption{不同测试族的存在域与识别能力。复指数、Cauchy 核和半轴 Laplace 核各有直接唯一性机制；
多项式矩在无界支集上需要紧支集、指数矩或 Carleman 一类增长条件。}
\label{fig:7-3-transform-determining-families}
\end{figure}

这张比较表给出选工具的原则。组合概率中的极限定理常能逐阶计算矩，适合使用
定理~\ref{theorem:7-3-3-4}；谱与预解式问题通常直接产生 Cauchy--Stieltjes 变换；卷积与独立
和则优先使用总存在的特征函数。无论选择哪一族测试，证明都要分成“先由紧性产生候选极限，
再由决定类识别候选”两步。弱收敛本身不能直接测试无界的 $x^k$，而全部矩存在也不能代替矩确定性。

\begin{remark}[三条不能省略的边界]
\begin{enumerate}
  \item 矩母函数可能除零点外处处为无穷，特征函数却对每个概率、每个实参数都存在；
  \item 弱收敛只控制有界连续测试。要传递第 $k$ 阶矩，必须另有更高阶矩给出一致可积尾界；
  \item 若极限矩问题不确定，可在例~\ref{ex:7-3-1-3} 的两个同矩分布之间交替。此时每一阶矩都
  恒定，分布却不收敛。
\end{enumerate}
\end{remark}

\begin{exercise}
设 $\mu,\nu$ 支持在同一紧集 $K\subset\R$ 且矩相同。把
例~\ref{ex:7-3-3-1} 的证明改写为对符号测度 $\mu-\nu$ 的全变差估计。
\end{exercise}
\begin{exercise}
不用完整 Stirling 公式，只用
$m_{2n}\leq(2n)^n$ 验证标准 Gaussian 的 Carleman 条件；再计算方差为 $\sigma^2$、均值为
$m$ 时条件仍成立。
\end{exercise}
\begin{exercise}
在定理~\ref{theorem:7-3-3-4} 中固定 $k$，写出取偶数 $q>k$ 后
$\{|X_n|^k\}$ 一致可积的完整估计，并说明只知道 $k$ 阶矩有界为什么不够。
\end{exercise}

本节从矩阵正性出发构造并识别测度，又用 Helly 与紧性控制测度列的极限。下一节将把
$f\mapsto\int f\,\dd\mu$ 反过来看成正线性泛函，证明 Daniell--Stone 与 Riesz--Markov 表示；
那时本节的有限构造会获得一个更统一、但逻辑上更晚的抽象解释。
